
\documentclass[12pt]{article}

% Layout.
\usepackage[top=1.2in, bottom=0.75in, left=1in, right=1in, headheight=1.0in, headsep=0pt]{geometry}

% Fonts.
\usepackage{mathptmx}
\usepackage[scaled=0.86]{helvet}
\renewcommand{\emph}[1]{\textsf{\textbf{#1}}}

% TiKZ.
\usepackage{tikz, pgfplots}
\usetikzlibrary{calc}
\pgfplotsset{my style/.append style={axis x line=middle, axis y line=middle, xlabel={$x$}, ylabel={$y$}}}
\pgfplotsset{compat=1.16}

% Misc packages.
\usepackage{amsmath,amssymb,latexsym,bm,array,multicol,enumitem}
\usepackage{graphicx}
\usepackage{xcolor}

% Commands to set various header/footer components.
\makeatletter
\def\doctitle#1{\gdef\@doctitle{#1}}
\doctitle{Use {\tt\textbackslash doctitle\{MY LABEL\}}.}
\def\docdate#1{\gdef\@docdate{#1}}
\docdate{Use {\tt\textbackslash docdate\{MY DATE\}}.}
\def\doccourse#1{\gdef\@doccourse{#1}}
\let\@doccourse\@empty
\def\docscoring#1{\gdef\@docscoring{#1}}
\let\@docscoring\@empty
\def\docversion#1{\gdef\@docversion{#1}}
\let\@docversion\@empty
\makeatother

% Headers and footers layout.
\makeatletter
\usepackage{fancyhdr}
\pagestyle{fancy}
\fancyhf{} % Clears all headers/footers.
\lhead{\emph{\@doctitle\hfill\@docdate} \medskip
\ifnum \value{page} > 1\relax\else\\
\emph{Name: \rule{3.5in}{1pt}\hfill \@docscoring}
\fi}

\rfoot{\emph{\@docversion}}
\lfoot{\emph{\@doccourse}}
\cfoot{\emph{\thepage}}
\renewcommand{\headrulewidth}{0pt}%
\makeatother

% Paragraph spacing
\parindent 0pt
\parskip 6pt plus 1pt

% A problem is a section-like command. Use \problem{5} for a problem worth 5 points.
\newcounter{probcount}
\newcounter{subprobcount}
\setcounter{probcount}{0}

\newcommand{\problem}[1]{%
\par
\addvspace{4pt}%
\setcounter{subprobcount}{0}%
\stepcounter{probcount}%
\makebox[0pt][r]{\emph{\arabic{probcount}.}\hskip1ex}\emph{[#1 points]}\hskip1ex}

\newcommand{\thesubproblem}{\emph{\alph{subprobcount}.}}

% Subproblems are an enumerate-like environment with a consistent
% numbering scheme.  Use \begin{subproblems}\item...\item...\end{subproblems}
\newenvironment{subproblems}{%
\begin{enumerate}[itemindent=8pt,leftmargin=0pt]%
\setcounter{enumi}{\value{subprobcount}}%
\renewcommand{\labelenumi}{\emph{\textsl{\alph{enumi})}} \,}%
}%
{%
\setcounter{subprobcount}{\value{enumi}}%
\end{enumerate}%
}

% Blanks for answers in normal and math mode.
\newcommand{\blank}[1]{\rule{#1}{0.75pt}}
\newcommand{\mblank}[1]{\underline{\hspace{#1}}}
\def\emptybox(#1,#2){\framebox{\parbox[c][#2]{#1}{\rule{0pt}{0pt}}}}

% Misc.
\renewcommand{\d}{\displaystyle}
\newcommand{\ds}{\displaystyle}
\newcommand{\threeopts}{{\small \hspace{-6mm} $\begin{matrix} \text{\textsc{converges}} \\ \text{\textsc{absolutely}} \end{matrix}$ \qquad\qquad $\begin{matrix} \text{\textsc{converges}} \\ \text{\textsc{conditionally}} \end{matrix}$ \qquad\qquad \textsc{diverges}} \bigskip}


\newcommand{\ba}{\mathbf{a}}
\newcommand{\bb}{\mathbf{b}}
\newcommand{\bc}{\mathbf{c}}
\newcommand{\bi}{\mathbf{i}}
\newcommand{\bj}{\mathbf{j}}
\newcommand{\bk}{\mathbf{k}}
\newcommand{\bu}{\mathbf{u}}
\newcommand{\bv}{\mathbf{v}}
\newcommand{\bw}{\mathbf{w}}

\newcommand{\ip}[2]{\mathrm{\left<#1,#2\right>}}
\newcommand{\vv}[2]{\mathrm{\left<#1,#2\right>}}
\newcommand{\vvv}[3]{\mathrm{\left<#1,#2,#3\right>}}



% Solution boxes.
\usepackage[most]{tcolorbox}
\usepgfplotslibrary{fillbetween}
\definecolor{solbg}{RGB}{225,240,252}
\definecolor{solframe}{RGB}{70,130,180}
\newtcolorbox{solution}[1][]{%
  breakable, enhanced,
  colback=solbg, colframe=solframe, boxrule=0.7pt, arc=3pt,
  left=6pt, right=6pt, top=5pt, bottom=5pt,
  before skip=6pt, after skip=6pt,
  fonttitle=\sffamily\bfseries\small, coltitle=white, title={Solution},
  #1}

\doctitle{Math 252: Quiz 4 Solutions}
\docdate{24 September, 2026}
\doccourse{}
\docversion{Solutions}
\docscoring{\fbox{{\LARGE \strut}\hspace{0.8in} / 25}}

\begin{document}
30 minutes maximum.  No aids (book, calculator, etc.) are permitted.
Show all work and use proper notation for full credit.  Answers should
be in reasonably-simplified form.  25 points possible.

\begin{enumerate}
\item (5 points per part) Evaluate the following integrals and
  antiderivatives.
  \begin{enumerate}
  \item $\ds \int_1^e x \ln(x) \, dx$

\begin{solution}
Integrate by parts with $u=\ln(x)$ and $dv=x\,dx$, so $du=\dfrac{1}{x}\,dx$ and
$v=\dfrac{x^2}{2}$:
\begin{align*}
  \int_1^e x\ln(x)\,dx
  &\;=\; \left[\frac{x^2}{2}\ln(x)\right]_1^e-\int_1^e \frac{x^2}{2}\cdot\frac{1}{x}\,dx \\[4pt]
  &\;=\; \frac{e^2}{2}-\frac{1}{2}\int_1^e x\,dx
  \;=\; \frac{e^2}{2}-\frac{1}{2}\left[\frac{x^2}{2}\right]_1^e \\[4pt]
  &\;=\; \frac{e^2}{2}-\frac{e^2}{4}+\frac{1}{4}.
\end{align*}
\[
\boxed{\ \int_1^e x\ln(x)\,dx=\frac{e^2+1}{4}\ }
\]
\end{solution}

  \item $\ds \int x^2 e^{3x} \, dx$

\begin{solution}
Integrate by parts with $u=x^2$ and $dv=e^{3x}\,dx$, so $du=2x\,dx$ and
$v=\dfrac{e^{3x}}{3}$:
\[
\int x^2e^{3x}\,dx=\frac{x^2e^{3x}}{3}-\frac{2}{3}\int xe^{3x}\,dx .
\]
Again with $u=x$ and $dv=e^{3x}\,dx$, so $du=dx$ and $v=\dfrac{e^{3x}}{3}$,
\[
\int xe^{3x}\,dx=\frac{xe^{3x}}{3}-\frac{1}{3}\int e^{3x}\,dx=\frac{xe^{3x}}{3}-\frac{e^{3x}}{9}.
\]
Therefore
\[
\int x^2e^{3x}\,dx=\frac{x^2e^{3x}}{3}-\frac{2xe^{3x}}{9}+\frac{2e^{3x}}{27}+C .
\]
\[
\boxed{\ \int x^2e^{3x}\,dx=\frac{e^{3x}}{27}\left(9x^2-6x+2\right)+C\ }
\]
\end{solution}

  \end{enumerate}

  \newpage
  
\item (5 points per part) Evaluate the following integrals and
  antiderivatives.
  \begin{enumerate}
  \item $\ds \int \sin^3(x) \cos^4(x) \, dx$

\begin{solution}
Split off one factor of $\sin(x)$, write $\sin^2(x)=1-\cos^2(x)$, and substitute
$u=\cos(x)$, $du=-\sin(x)\,dx$:
\begin{align*}
  \int \sin^3(x)\cos^4(x)\,dx
  &\;=\; \int\left(1-\cos^2(x)\right)\cos^4(x)\sin(x)\,dx \\[4pt]
  &\;=\; -\int\left(1-u^2\right)u^4\,du
  \;=\; \int\left(u^6-u^4\right)du
  \;=\; \frac{u^7}{7}-\frac{u^5}{5}+C .
\end{align*}
\[
\boxed{\ \int \sin^3(x)\cos^4(x)\,dx=\frac{\cos^7(x)}{7}-\frac{\cos^5(x)}{5}+C\ }
\]
\end{solution}

  \item $\ds \int \tan^6(x) \sec^4(x) \, dx$

\begin{solution}
Split off one factor of $\sec^2(x)$, write $\sec^2(x)=1+\tan^2(x)$, and substitute
$u=\tan(x)$, $du=\sec^2(x)\,dx$:
\begin{align*}
  \int \tan^6(x)\sec^4(x)\,dx
  &\;=\; \int \tan^6(x)\left(1+\tan^2(x)\right)\sec^2(x)\,dx \\[4pt]
  &\;=\; \int\left(u^6+u^8\right)du
  \;=\; \frac{u^7}{7}+\frac{u^9}{9}+C .
\end{align*}
\[
\boxed{\ \int \tan^6(x)\sec^4(x)\,dx=\frac{\tan^7(x)}{7}+\frac{\tan^9(x)}{9}+C\ }
\]
\end{solution}

  \item $\ds \int_{-\pi / 2}^{\pi / 2} \cos^3(x) \, dx$

\begin{solution}
Split off one factor of $\cos(x)$ and write $\cos^2(x)=1-\sin^2(x)$.  With
$u=\sin(x)$, $du=\cos(x)\,dx$, the limits become $u=-1$ and $u=1$:
\[
\int_{-\pi/2}^{\pi/2}\cos^3(x)\,dx=\int_{-1}^{1}\left(1-u^2\right)du
=\left[u-\frac{u^3}{3}\right]_{-1}^{1}=\frac{2}{3}-\left(-\frac{2}{3}\right).
\]
\[
\boxed{\ \int_{-\pi/2}^{\pi/2}\cos^3(x)\,dx=\frac{4}{3}\ }
\]
\end{solution}

  \end{enumerate}
\end{enumerate}


\end{document}

%%% Local Variables:
%%% mode: latex
%%% TeX-master: t
%%% End:
